% Proposed replacement for the abstract and Section 1 of main.tex (lines 9--15 and 19--168 of the fourth draft as of 03:31 UTC, 24 November 2026).
% Theorems 1.1 and 1.2 are verbatim. In Theorem 1.3 only the sentences introducing x, y and the gate are
% reordered; the hypotheses and the conclusion are unchanged. The list "What is not claimed" is verbatim.
% Labels kept: s:intro, t:crit, t:inh, e:hyp, e:inh, t:ret, e:ret, s:not.

\begin{abstract}
Suppose a Yang--Mills theory is known only through its correlation functions. How could a mass gap be
recognised in it, and what would let a gap survive when a lattice is refined? For the first question we prove
a spectral criterion. For the second we prove an inequality that reduces one refinement step to three
estimates, and we prove one of the three for $\SU(2)$. The other two are open, and so is a further condition
needed to chain the steps. No mass gap is proved.
\end{abstract}

\section{Introduction}\label{s:intro}

Suppose that a quantum Yang--Mills theory in four dimensions has been constructed, as a limit of lattice
gauge theories, and that one is handed its correlation functions. Two questions come to mind.
\begin{enumerate}
\item How would one see that the theory has a mass gap?
\item If the lattice theories have gaps in a fixed physical box, what would allow a gap to survive as the
mesh is halved again and again?
\end{enumerate}
By a \emph{mass gap} we mean a number $\Delta>0$ such that the Hamiltonian $H\ge0$ has no spectrum in
$(0,\Delta)$ and has some spectrum above it. The lattice terms used in this introduction (link, plaquette,
mesh, fixed physical box) are defined in Section~\ref{s:set}.

This paper does not prove a mass gap. It records what we can prove about the two questions, states what is
open, and reports measurements of the quantities that the open statements concern.

The first question has recently become more concrete. A construction of a continuum theory, as a scaling
limit of Wilson lattice gauge theory for every simple compact gauge group, was submitted in formal form in
October 2026~\cite{OAI}. Its assessors judge that it satisfies the Wightman axioms, and state that it proves
neither a mass gap nor uniqueness of the infinite-volume limit~\cite[p.~1]{YMG}. The problem of Jaffe and
Witten~\cite[\S4]{JW} asks for existence and a gap. Whether that construction is accepted as existence is not
for us to say, and it has further conditions (Section~\ref{s:ext}); if it is, a gap is what would still be
missing.

\subsection{The first question}
Here the answer is an elementary fact about self-adjoint operators, worth stating because it fits a theory
known through correlation functions. If $H\ge0$ is the Hamiltonian and $v$ a vector, the numbers
$\langle v,e^{-n\tau H}v\rangle$ are correlation functions at Euclidean time separation $n\tau$. Their
successive differences $d_{v,n}$ are non-negative, receive no contribution from the vacuum, and decay at a
rate set by the lowest positive energy present in $v$. Theorem~\ref{t:crit} says that $H$ has no spectrum in
$(0,\Delta)$ exactly when these differences decay at one common rate for a family of vectors with dense
linear span.

\subsection{The second question: one route}
Draw a square grid, and over it a second grid with squares twice as wide. A lattice gauge field puts a group
element on each edge of the grid; the edges are called \emph{links}. A field on the fine grid gives one on
the coarse grid: multiply the two fine links under each coarse link. Any function of the fine field at one
time is then the sum of its best guess from the coarse field and a remainder, the detail that the coarse
field does not see.

Suppose three things.
\begin{enumerate}
\item[(a)] The guess decays in time as on the coarse lattice.
\item[(b)] Every function, whatever its guess, loses within two time steps a fixed fraction of its detail,
except for detail carried by an exceptional set of coarse fields.
\item[(c)] The dynamics seldom returns to that exceptional set.
\end{enumerate}
Then the fine lattice keeps the coarse decay up to two losses, which must both be small: the return to
the exceptional set, and a factor that is close to one when the detail is lost in few steps compared with
the time considered. Theorem~\ref{t:inh} makes this precise.

The exceptions are allowed for a reason. A natural way to prove~(b) goes through local updates of the fine
links, and that route meets slow trial directions on selected coarse fields (Section~\ref{s:local}).

Where this stands: Theorem~\ref{t:ret} proves a bound of the kind~(c) for $\SU(2)$. Statements (b) and~(a)
are open (Hypotheses~\ref{h:mix} and~\ref{h:coarse}). Chaining the refinement steps needs, in addition, a
bound on the accumulated losses, which is also open (Hypothesis~\ref{h:budget}).

\subsection{Main results}
\begin{theorem}[Spectral criterion]\label{t:crit}
Let $H\ge0$ be a self-adjoint operator on a Hilbert space $\cH$, let $\tau>0$, and let $D\subset\cH$ have dense
linear span. For $v\in\cH$ and $n\ge0$ put
\[
d_{v,n}=\big\langle v,\,e^{-n\tau H}\big(I-e^{-\tau H}\big)v\big\rangle\ \ge0 .
\]
Let $0<r<1$ and $\Delta=-\tau^{-1}\log r$. The following are equivalent.
\begin{enumerate}
\item $H$ has no spectrum in $(0,\Delta)$.
\item $\limsup_{n\to\infty}d_{v,n}^{1/n}\le r$ for every $v\in D$.
\end{enumerate}
\end{theorem}

Only the rate $r$ is common in (2); constants and the onset of the decay may depend on $v$. The criterion is
satisfied trivially when $H=0$, so by itself it says nothing about whether there is any positive energy; the
problem of~\cite{JW} also requires the mass to be finite.

\begin{theorem}[Refinement with an exceptional set]\label{t:inh}
Let $A$ be a self-adjoint operator on a non-zero Hilbert space $\mathcal K$ with $0\le A\le I$. Let $P$ be an
orthogonal projection, $Q=I-P$, and let $B$ be an orthogonal projection with range inside the range of $Q$;
put $G=Q-B$. Let $m\ge j\ge1$ be integers, $\gamma>0$, $\kappa\ge0$ and $b\in\R$, and suppose
\begin{equation}\label{e:hyp}
\langle f,(I-A^{2})f\rangle\ge\gamma\|Gf\|^2\ \ (f\in\mathcal K),\qquad
\|BA^{j}B\|\le\kappa,\qquad \|PA^{m}P\|\le1-b .
\end{equation}
Then
\begin{equation}\label{e:inh}
1-\|A^m\|\ \ge\ \frac{(b-\kappa)_+}{1+2/(m\gamma)} .
\end{equation}
\end{theorem}

Here $(s)_+=\max\{s,0\}$. The space splits into a coarse part $P$, an exceptional part $B$ and a good part
$G$, and the three hypotheses are (b), (c) and (a) above, in that order:
\begin{itemize}
\item a function loses, within two steps, at least the fraction $\gamma$ of the squared norm of its good
part;
\item the exceptional part overlaps with itself after $j$ steps by at most $\kappa$;
\item the coarse part alone already has deficit $b$ after $m$ steps.
\end{itemize}
The conclusion is that the whole space has deficit at least $b-\kappa$, divided by a factor that tends to $1$
when $m\gamma\to\infty$.

In the application $A$ is one step of a transfer semigroup. If $A=e^{-aH}$ with $H\ge0$ and mesh $a$, then
$A^m=e^{-tH}$ for $t=ma$, the left side of~\eqref{e:inh} is the \emph{deficit} $d=1-e^{-t\inf\sigma(H)}$ at
physical time $t$, and the denominator is $1+2a/(\gamma t)$. No logarithm of $A$ is needed, and $A$ may have
a kernel. In detail (terms are defined in Section~\ref{s:set}; symbols are listed in Table~\ref{tab:sym}),
$\mathcal K$ is the space of gauge-invariant functions of one time slice of the fine lattice with vacuum
expectation zero, $A$ is one step of the transfer semigroup in the vacuum representation, $P$ is the
conditional expectation given the coarse links, and $B$ is the part of $Q$ supported on an exceptional set of
coarse fields. A physical time $t$ corresponds to $m=t/a$ steps, so $m$ doubles when the mesh is halved. On
this space $1-\|A^m\|>0$ is a gap of that lattice.

\begin{theorem}[Return to a coarse plaquette gate]\label{t:ret}
Consider the Wilson measure with weight $\exp\big(\beta\sum_p\tfrac12\Tr U_p\big)$, $\beta\ge0$, for the gauge
group $\SU(2)$ on a periodic lattice with even spatial side $N\ge6$ and $T\ge2$ time slices. Fix a spatial
plaquette of side two (a coarse plaquette) and $0\le\theta\le\pi$. Let $x$ and $y$ be the spatial links of
two time slices that are distinct modulo $T$. For the spatial links of a time slice let $\Gamma_\theta$ be the
event that the holonomy around that plaquette (the ordered product of its eight fine links) lies within round
distance $\theta$ of $-I$ (the angle between the two as points of $S^3$); we call $\Gamma_\theta$ a
\emph{gate}. Then for every complex square-integrable function $g$ of a slice, with
$h=\mathbf 1_{\Gamma_\theta}g$,
\begin{equation}\label{e:ret}
\big|\E\big[\overline{h(x)}\,h(y)\big]\big|\ \le\ p_{6\beta}(\theta)\,\E\big[|h(x)|^2\big],
\qquad
p_k(\theta)=\frac{\int_0^\theta\sin^2\!r\,e^{k(\cos r-1)}\,dr}{\int_0^\pi\sin^2\!r\,e^{k(\cos r-1)}\,dr},
\end{equation}
and $p_k(\theta)\le\min\{1,\ \tfrac{\pi^2\sqrt e}{4}(1+k)^{3/2}\theta^3\}$.
\end{theorem}

In words: whatever the field is at one time, the conditional probability that the chosen plaquette is within
$\theta$ of $-I$ at another time is at most $p_{6\beta}(\theta)$, which is of order $\theta^3$. Three remarks
on its reach.
\begin{itemize}
\item The order $\theta^3$ is for fixed $\beta$; the bound is small only for $\theta$ well below
$(6\beta)^{-1/2}$.
\item The bound does not depend on the spatial volume, on the time separation or on the rest of the field.
\item The gate is centred at $-I$ because the smallest local quotients we measured are on a background with
three plaquettes at $-I$ (Section~\ref{s:local}: two backgrounds, one lattice, one coupling); this is a
choice suggested by those trials, not a result.
\end{itemize}
Corollary~\ref{c:vac} carries the bound to the vacuum on an infinite time axis, where it becomes the second
inequality of~\eqref{e:hyp}.

Two further results are proved. Proposition~\ref{p:three} in Section~\ref{s:crit} is a form of the criterion
that uses a single time separation. Proposition~\ref{p:haar} in Section~\ref{s:haar} gives an explicit lower
bound for the conditional variance of a single link; it bounds one term of the auxiliary statement of
Section~\ref{s:arch} and is not an input of Theorem~\ref{t:inh}.

\subsection{What is not claimed}\label{s:not}
\begin{enumerate}
\item No mass gap is proved in infinite volume or in the continuum, and no lower bound uniform along a
refinement, for any gauge group. (At a fixed finite lattice a gap exists for elementary reasons and is not
the issue.)
\item No continuum theory is constructed. We have not examined the construction~\cite{OAI}; we rely on its
assessment~\cite{YMG} for what it does and does not contain, and we have not shown that any object of this
paper belongs to it (Section~\ref{s:ext}).
\item Theorem~\ref{t:inh} is conditional. Of its three hypotheses the second is proved for $\SU(2)$
(Corollary~\ref{c:vac}). The first and the third are open (Hypotheses~\ref{h:mix} and~\ref{h:coarse}), and so
is the budget that chaining the steps needs (Hypothesis~\ref{h:budget}); none is proved on any trajectory of
matched couplings. Even with all of them the conclusion concerns a fixed physical box and is not a gap in
infinite volume.
\item Theorem~\ref{t:ret} and all measurements are for $\SU(2)$. The gate around $-I$ and the constant $p_{6\beta}(\theta)$
are special to it; that constant is the largest value over the conditional laws that one link can have, and
is not shown to be attained on any Wilson background. For other groups only a cruder return bound is available (Remark~\ref{r:ret}(d)), and
whether a useful exceptional set exists is open.
\item The measurements of Section~\ref{s:evid} are Monte Carlo estimates on lattices of side $5$ to $10$,
one coupling each, with statistical errors. The exact value of a trial function limits a constant from one
side; an estimate of it is not a bound. No measurement is evidence for or against a hypothesis, and none is
used in a proof.
\end{enumerate}

How the results were obtained, which parts are machine-checked, and the fact that the paper has not been
peer reviewed are set out in Section~\ref{s:disc}.
